% BEGIN REV09_GERMEN_DODECAFASICO
\section{Germen dodecafásico y recuperación de las ramas unimodales}
\label{corpus9:iv:germen}

La rueda orientada \(U_{12}^{\mathrm{sign}}\), recibida del retorno
TPK, determina en el sector uniforme \(\eta=0\) los pesos
\(w_m=1/12\). Fijamos el funcional de fase y sus representantes
reales ordenados
\begin{equation}
 \phi_m=\pi_{\mathrm{HMT}}\frac{3m-1}{18},
 \qquad m=1,\ldots,12.
 \label{corpus9:iv:germen-fases}
\end{equation}
El levantamiento real de las fases forma parte del lector:
los momentos siguientes se calculan sobre esos representantes.
La construcción procede de la rueda, su orientación, sus pesos
y su funcional; una modificación de cualquiera de esos datos
modifica la colección resultante.

\subsection{Normalización y perfil analítico}
Sean \(a_m=(3m-1)/18\) y
\(D_0=\frac1{12}\sum_{m=1}^{12}a_m^2\). Las sumas enteras
\[
 \sum_{m=1}^{12}(3m-1)^2=5394=6\cdot899,\qquad
 \sum_{m=1}^{12}(3m-1)^4=4294614
\]
dan \(D_0=899/648\). La condición de normalización cuadrática
\(s_0^2\pi_{\mathrm{HMT}}^2D_0=2\) fija
\[
 s_0=\frac{36}{\pi_{\mathrm{HMT}}\sqrt{899}},
 \qquad k_m=s_0\phi_m=\frac{2(3m-1)}{\sqrt{899}}.
\]
La escala angular se cancela en esta normalización antes de
iterar el perfil
\begin{equation}
 u_0(x)=1-\frac1{12}\sum_{m=1}^{12}\cos(k_mx),
 \qquad f_\lambda(x)=1-\lambda u_0(x).
 \label{corpus9:iv:germen-perfil}
\end{equation}

\begin{proposition}[Representación exacta y criticidad]
\label{corpus9:iv:germen-analitico}
El perfil \(u_0\) es una función entera par y satisface
\[
 u_0(x)=x^2-\frac{715769}{2424603}x^4+O(x^6).
\]
Con \(y=x/\sqrt{899}\), su expresión cerrada es
\begin{equation}
 u_0(x)
  =1-\frac{\cos(37y)\sin(36y)}{12\sin(3y)}
  =1-\frac{\sin(73y)-\sin y}{24\sin(3y)},
 \label{corpus9:iv:germen-cerrado}
\end{equation}
donde los cocientes se prolongan en sus singularidades removibles.
Para \(0<\lambda\leq2/u_0(1)\), \(f_\lambda\) es una
autoaplicación real analítica y par de \([-1,1]\), con un
único punto crítico cuadrático en cero y Schwarziano negativo
fuera de ese punto.
\end{proposition}
\begin{proof}
La suma finita de cosenos es entera y par. Su desarrollo de
Taylor utiliza
\[
 \frac1{12}\sum_m k_m^2=2,\qquad
 \frac1{24\cdot12}\sum_m k_m^4=\frac{715769}{2424603},
\]
obtenidos de las dos sumas enteras anteriores. Para la forma
cerrada, la progresión de argumentos es \((6m-2)y\). La suma
geométrica de los exponenciales, seguida de la parte real, da
\[
 \sum_{m=1}^{12}\cos((6m-2)y)
   =\frac{\sin(36y)}{\sin(3y)}\cos(37y).
\]
La identidad \(2\cos(37y)\sin(36y)=\sin(73y)-\sin y\)
prueba \eqref{corpus9:iv:germen-cerrado}. La suma finita define
su prolongación en todos los ceros del denominador.

Para \(0<x\leq1\), todos los argumentos satisfacen
\(0<k_mx\leq70/\sqrt{899}<\pi_{\mathrm{HMT}}\). Por tanto,
\[
 u_0'(x)=\frac1{12}\sum_m k_m\sin(k_mx)>0,\qquad
 u_0'''(x)=-\frac1{12}\sum_m k_m^3\sin(k_mx)<0.
\]
La paridad proporciona los signos opuestos en el intervalo
negativo y prueba la unicidad del punto crítico. Además,
\(f_\lambda''(0)=-2\lambda\ne0\). Como
\(0\leq u_0(x)\leq u_0(1)\) en \([-1,1]\), el intervalo
indicado para \(\lambda\) da
\(-1\leq f_\lambda(x)\leq1\). Finalmente,
\[
 S(f_\lambda)
  =\frac{u_0'''}{u_0'}
   -\frac32\left(\frac{u_0''}{u_0'}\right)^2<0
 \qquad(0<|x|\leq1).
\]
El primer cociente es negativo en ambos lados del origen,
y el término restante es no positivo.
\end{proof}

\subsection{Inversión por rama y memoria del itinerario}
\begin{proposition}[Recuperación del plegamiento]
\label{corpus9:iv:germen-inversa}
Sea \(\lambda>0\) y
\(v=(u_0|_{[0,1]})^{-1}\). Para
\(y\in[1-\lambda u_0(1),1]\), la preimagen en una rama es
\begin{equation}
 x=\sigma\,v\!\left(\frac{1-y}{\lambda}\right),
 \qquad \sigma=\operatorname{sgn}x.
 \label{corpus9:iv:germen-rama}
\end{equation}
En \(y=1\), la preimagen es cero y se fija \(\sigma=0\).
\end{proposition}
\begin{proof}
La continuidad y la monotonía estricta de \(u_0\) en
\([0,1]\) proporcionan la inversa \(v\).
La ecuación \(y=1-\lambda u_0(x)\), junto con la paridad,
determina \(|x|=v((1-y)/\lambda)\); el signo determina la
orientación. Para \(y=1\), la unicidad del mínimo de \(u_0\)
da \(x=0\). Fuera del punto crítico, la inversa es diferenciable.
Su derivada se singulariza al llegar al origen, donde
\(u_0(x)=x^2+O(x^4)\).
\end{proof}

Para una distribución finita sobre pares \(\pm x_i\), con
\(x_i>0\) distintos, sean \(p_i^\pm\) sus probabilidades,
\(p_i=p_i^++p_i^-\), \(Y=f_\lambda(X)\) y
\(\Sigma=\operatorname{sgn}X\). Se admite también una masa
en cero, cuyo término de incertidumbre de rama es nulo. Con
logaritmos naturales,
\begin{align}
 H(X\mid Y,\Sigma)&=0,\\
 H(X\mid Y)&=H(\Sigma\mid Y)
   =\sum_{i:p_i>0}p_i\,h(p_i^+/p_i),\\
 h(p)&=-p\log p-(1-p)\log(1-p).
 \label{corpus9:iv:germen-entropia}
\end{align}
La prueba consiste en agrupar la distribución por las fibras
de \(f_\lambda\). La monotonía distingue los módulos \(x_i\),
y \eqref{corpus9:iv:germen-rama} distingue los dos signos
dentro de cada fibra. Con probabilidades iguales en cada par
y sin masa en cero, la incertidumbre vale \(\log2\).

Si \(0<\lambda\leq2/u_0(1)\), la segunda iteración está
definida sobre el mismo intervalo y satisface
\[
 H(X\mid f_\lambda^2(X))
 =H(X\mid f_\lambda(X))
  +H(f_\lambda(X)\mid f_\lambda^2(X)).
\]
En efecto, ambas aplicaciones intermedias son deterministas,
y la regla de la cadena para la entropía de
\((X,f_\lambda(X))\), condicionada a \(f_\lambda^2(X)\),
da la identidad. La recuperación sucesiva utiliza los signos
de cada paso. Para \(\lambda=0\), la aplicación es constante
y requiere conservar además el módulo de la entrada.

\subsection{Ámbito del retorno reescalado}
Para \(a=-f(1)\ne0\), sea \(h_a(x)=-ax\). El operador
\[
 \mathcal R(f)=h_a^{-1}\circ f^2\circ h_a
\]
es un retorno reescalado sobre \([-1,1]\) cuando el intervalo
\(J=h_a([-1,1])\) pertenece al dominio de \(f^2\) y cumple
\(f^2(J)\subseteq J\). La escala \(h_a\) es biyectiva;
el retorno de dos pasos conserva su interpretación con el
itinerario y la aplicación \(f\). La reconstrucción de una
raíz iterativa a partir de \(\mathcal R(f)\) y \(a\) constituye
una condición adicional. La proposición
\ref{corpus9:iv:germen-analitico} establece el germen y su
clase crítica. La existencia y unicidad de un punto fijo del
renormalizador, su descomposición espectral estable e inestable
y la transversalidad de esta colección son las hipótesis
adicionales del paso a un límite universal de renormalización.
% END REV09_GERMEN_DODECAFASICO
